
\begin{definition} [Multifold sums of powers]
  \label{def:multifold-sums-of-powers-recurrence}
  For non-negative integers $r,n,m$
  \begin{align*}
    \KnuthRFoldSum{0}{n}{m}   &= n^m, \\
    \KnuthRFoldSum{1}{n}{m}   &= \KnuthRFoldSum{0}{1}{m}
    + \KnuthRFoldSum{0}{2}{m} + \cdots + \KnuthRFoldSum{0}{n}{m}, \\
    \KnuthRFoldSum{r+1}{n}{m} &= \KnuthRFoldSum{r}{1}{m}
    + \KnuthRFoldSum{r}{2}{m} + \cdots + \KnuthRFoldSum{r}{n}{m}.
  \end{align*}
\end{definition}
\begin{definition}[Central factorials]
  \label{def:central-factorials}
  For integers $n,k$
  \begin{align*}
    \centralFactorial{n}{k} =
    \begin{cases}
      0, \quad & \mathrm{if \; } k < 0, \\
      1, \quad & \mathrm{if \; } k = 0, \\
      n \left( n + \frac{k}{2} -1 \right)\left( n + \frac{k}{2} -2 \right)
      \cdots \left( n - \frac{k}{2} +1 \right)
      = n \prod_{j=1}^{k-1} \left( n + \frac{k}{2} -j \right), \quad & \mathrm{if \; } k>0.
    \end{cases}
  \end{align*}
\end{definition}
\begin{proposition}[Central Newton's formula]
  For arbitrary integers $x$, and $t$
  \label{prop:newtons-interpolation-formula-in-central-differences}
  \begin{align*}
    f(x)
    = \sum_{k=0}^{\infty} \frac{\centralFactorial{(x-t)}{k}}{k!} \delta^{k} f(t),
  \end{align*}
  where $\delta^{k} f(t)$ denotes central finite difference of $f(t)$,
  defined by
  \begin{align*}
    \delta^{k} f(t) = \tsum_{j=0}^{k} (-1)^j \tbinom{k}{j} f(t+\tfrac{k}{2} - j).
  \end{align*}
  \begin{proof}
    See~\cite[p. 462]{butzer1989central}.
  \end{proof}
\end{proposition}
\begin{lemma}[Central falling factorials]
  \label{lem:central-factorials-in-terms-of-falling}
  For integers $n,k$
  \begin{align*}
    \centralFactorial{n}{k} =
    \begin{cases}
      0, \quad & \mathrm{if \; } k<0, \\
      1, \quad & \mathrm{if \; } k=0, \\
      n \fallingFactorial{n+\frac{k}{2}-1}{k-1}, \quad & \mathrm{if \; } k>0.
    \end{cases}
  \end{align*}
  where
  $\fallingFactorial{n+\frac{k}{2}-1}{k-1}
  = \prod_{j=1}^{k-1} \left( n + \frac{k}{2} -j \right)$
  are falling factorials.
\end{lemma}
% \begin{lemma}[Binomial form of central factorials]
%   \label{lem:binomial-form-of-central-factorials}
%   For integers $n$ and $k \geq 1$,
%   \begin{align*}
%     \frac{\centralFactorial{n}{k}}{k!}
%     = \frac{n}{k} \binom{n+\frac{k}{2}-1}{k-1}.
%   \end{align*}
%   \begin{proof}
%     We have
%     \begin{align*}
%       \frac{\centralFactorial{n}{k}}{k!}
%       &= \frac{n}{k!}\fallingFactorial{n+\frac{k}{2}-1}{k-1}
%       = \frac{n}{k (k-1)!} \fallingFactorial{n+\frac{k}{2}-1}{k-1}
%       = \frac{n}{k} \binom{n+\frac{k}{2}-1}{k-1},
%     \end{align*}
%     because of the identity in falling factorials
%     $\frac{\fallingFactorial{x}{n}}{n!} = \binom{x}{n}$,
%     and Lemma~\eqref{lem:central-factorials-in-terms-of-falling}.
%   \end{proof}
% \end{lemma}
\begin{lemma}[Halved central factorials]
  \label{lem:halved-central-factorials}
  For integers $n,k\geq 0$,
  \begin{align*}
    \frac{\centralFactorial{n}{k}}{k!}
    = \frac{1}{2} \binom{n+\frac{k}{2}}{k} + \frac{1}{2} \binom{n+\frac{k}{2}-1}{k}.
  \end{align*}
  \begin{proof}
    We have $\binom{a}{k} = \frac{a}{k} \binom{a-1}{k-1}$.
    Let $a=n + \frac{k}{2}$, then
    \begin{align*}
      \tbinom{n + \frac{k}{2}}{k} = \tfrac{n + \tfrac{k}{2}}{k} \tbinom{n+\frac{k}{2}-1}{k-1}
      = \tfrac{n}{k} \tbinom{n+\frac{k}{2}-1}{k-1} + \tfrac{1}{2} \tbinom{n+\frac{k}{2}-1}{k-1}.
    \end{align*}
    Thus,
    \begin{align*}
      \tfrac{n}{k} \tbinom{n+\frac{k}{2}-1}{k-1} = \tbinom{n + \frac{k}{2}}{k} - \tfrac{1}{2} \tbinom{n+\frac{k}{2}-1}{k-1}.
    \end{align*}
    By moving under common denominator, and applying binomial recurrence yields
    \begin{align*}
      \tfrac{n}{k} \tbinom{n+\frac{k}{2}-1}{k-1}
      = \tfrac{1}{2}
      \left[
        2 \tbinom{n + \frac{k}{2}}{k} - \tbinom{n+\frac{k}{2}-1}{k-1}
      \right]
      = \tfrac{1}{2}
      \left[
        \tbinom{n + \frac{k}{2}}{k} + \tbinom{n + \frac{k}{2}-1}{k} + \tbinom{n + \frac{k}{2}-1}{k-1} - \tbinom{n+\frac{k}{2}-1}{k-1}
      \right].
    \end{align*}
    Thus, the statement follows.
  \end{proof}
\end{lemma}
\begin{proposition}[Central Newton's formula for powers]
  \label{prop:newtons-formula-for-powers}
  For integers $n,m \geq 0$, and an arbitrary integer $t$,
  \begin{align*}
    n^m = \sum_{k=0}^{m} \frac{\centralFactorial{(n-t)}{k}}{k!} \delta^{k} t^m.
  \end{align*}
\end{proposition}
\begin{proposition}[Halved central factorial powers]
  \label{prop:halved-central-factorial-powers}
  For integers $n,m \geq 0$, and an arbitrary integer $t$,
  \begin{align*}
    n^m = \frac{1}{2} \sum_{k=0}^{m}
    \left[
      \binom{n-t+\frac{k}{2}}{k} + \binom{n-t+\frac{k}{2}-1}{k}
    \right] \delta^{k} t^m
  \end{align*}
\end{proposition}
% \begin{proposition}[Generalized hockey-stick identity]
%   \label{prop:generalized-hockey-stick-identity}
%   For integers $a,b$ and $j$,
%   \begin{align*}
%     \sum_{k=a}^{b} \binom{k}{j} = \binom{b+1}{j+1} - \binom{a}{j+1}.
%   \end{align*}
%   \begin{proof}
%     We have,
%     $\sum_{k=a}^{b} \binom{k}{j} = \binom{a}{j} + \binom{a+1}{j} + \cdots + \binom{b}{j}$,
%     which means that,
%     $\sum_{k=a}^{b} \binom{k}{j} =
%     \left( \sum_{k=0}^{b} \binom{k}{j} \right) - \left( \sum_{k=0}^{a-1} \binom{k}{j} \right)$.
%     By hockey stick identity $\sum_{k=0}^{n} \binom{k}{j} = \binom{n+1}{j+1}$ yields,
%     \begin{align*}
%       \tsum_{k=a}^{b} \tbinom{k}{j}
%       = \left( \tsum_{k=0}^{b} \tbinom{k}{j} \right) - \left( \tsum_{k=0}^{a-1} \tbinom{k}{j} \right)
%       = \tbinom{b+1}{j+1} - \tbinom{a}{j+1}.
%     \end{align*}
%     This completes the proof.
%   \end{proof}
% \end{proposition}
% \begin{lemma}[Centered hockey-stick identity]
%   For integers $n\geq1$, and arbitrary integers $k,r$
%   \begin{align*}
%     \sum_{j=1}^{n} \binom{j + \frac{k}{2} +r}{k}
%     = \binom{n + \frac{k}{2} + r +1}{k+1} - \binom{\frac{k}{2} + r +1}{k+1}.
%   \end{align*}
%   \begin{proof}
%     By setting $j=1 + \tfrac{k}{2} +r$, and $b=n+ \tfrac{k}{2} +r$ yields
%     \begin{align*}
%       \tsum_{j=1}^{n} \tbinom{j + \tfrac{k}{2} +r}{k} =
%       \tsum_{j=1 + \tfrac{k}{2} +r}^{n+ \tfrac{k}{2} +r} \tbinom{j}{k}.
%     \end{align*}
%     By Generalized hockey-stick identity~\eqref{prop:generalized-hockey-stick-identity},
%     \begin{align*}
%       \tsum_{j=1 + \tfrac{k}{2} +r}^{n+ \tfrac{k}{2} +r} \tbinom{j}{k}
%       = \tbinom{n + \tfrac{k}{2} + r +1}{k+1} - \tbinom{\tfrac{k}{2} + r +1}{k+1}.
%     \end{align*}
%     Thus, the claim follows.
%   \end{proof}
% \end{lemma}
